% Generated draft: requires mathematical and source review.
\documentclass[11pt]{article}
\usepackage[margin=25mm]{geometry}
\usepackage{fontspec}
\setmainfont{Latin Modern Roman}
\usepackage{amsmath,amssymb,mathtools,unicode-math}
\setmathfont{Latin Modern Math}
\usepackage{enumitem,needspace,xurl,hyperref}
\setlength{\parindent}{0pt}
\setlength{\parskip}{5pt}
\setlength{\emergencystretch}{3em}
\newcommand{\N}{\mathbb{N}}
\newcommand{\Z}{\mathbb{Z}}
\newcommand{\Q}{\mathbb{Q}}
\newcommand{\R}{\mathbb{R}}
\newcommand{\C}{\mathbb{C}}
\newcommand{\F}{\mathbb{F}}
\newcommand{\A}{\mathbb{A}}
\newcommand{\PP}{\mathbb{P}}
\newcommand{\E}{\mathbb{E}}
\newcommand{\eps}{\varepsilon}
\newcommand{\dd}{\,\mathrm{d}}
\begin{document}
\section*{UnsolvedMath ID: 701717}
\textbf{Source problem code:} catalogue entry
\subsection{Definitions and mathematical statement}
A group is finitely generated if it is generated by a finite subset. A group \(G\) is virtually free if it contains a free subgroup \(F\leq G\) of finite index, meaning \([G:F]<\infty\). A subgroup \(H\leq K\) is a retract of \(K\) if there is a group homomorphism \(r:K\to H\) satisfying \(r(h)=h\) for every \(h\in H\). A subgroup \(H\leq G\) is a virtual retract of \(G\) if there exists a subgroup \(K\) such that \(H\leq K\leq G\), \([G:K]<\infty\), and \(H\) is a retract of \(K\).

The question proposed by A. Minasyan in Kourovka Notebook Problem 20.59 (2022) is: for every finitely generated virtually free group \(G\) and every finitely generated subgroup \(H\leq G\), do there exist a subgroup \(K\) and a group homomorphism \(r\) such that
\[
H\leq K\leq G,\qquad [G:K]<\infty,\qquad r:K\to H,\qquad \forall h\in H,\ r(h)=h?
\]
The supplied formulation imposes no normality requirement on \(K\).
\subsection{Short English statement}
Must every finitely generated subgroup of a finitely generated virtually free group be a retract of some finite-index subgroup containing it?
\subsection{Sources}
\begin{itemize}
\item Local UnsolvedMath registry, supplied complete source record: permanent ID \(701717\), exact source code \texttt{catalogue entry}, title \emph{Kourovka Notebook Problem 20.59}. The supplied registry labels the problem open; this label has not been independently verified.
\item \emph{The Kourovka Notebook (Unsolved Problems in Group Theory)}, 20th issue (2022), Problem 20.59, proposed by A. Minasyan. The supplied background also identifies the 21st issue edition (2026), \texttt{arXiv:1401.0300v48}. The registry describes the Issues 1--21 collection as edited by E. I. Khukhro and V. D. Mazurov.
\item A. Minasyan, \emph{Int. Math. Res. Notes}, \textbf{2021}, no. 17 (2021), 13434--13477. The supplied statement cites this reference for motivation and partial results; no specific partial-result statement is supplied.
\end{itemize}
\end{document}
